\subsection{Component-specific Gaussian precisions}
\label{sec:heterogeneous}

The sampled-time criterion also applies when the driving components
have different covariance matrices. The decision rule is still a
fixed affine readout. The class scores of the driving mixture, if
used as a readout, may define a different partition with curved
boundaries.

Let $a_i\in\mathbb R^d$, $w_i>0$, and let each $P_i$ be a symmetric
positive definite precision matrix. Write
\[
 \ell_i(x)=-\tfrac12(x-a_i)^\top P_i(x-a_i),\qquad
 U(x)=\log\sum_i w_i e^{\ell_i(x)},\qquad F=\nabla U.
\]
Component normalization constants can be included in the $w_i$.
Let $H=o+E$ be a facet hyperplane, with unit normal $n$,
$E=n^\perp$, and orthogonal projection $\Pi$ onto $E$. On $H$
the log contributions and their normal derivatives take the form
\[
 \begin{aligned}
 \ell_i(o+y)&=q_i(y)+c_i, &
 q_i(y)&=-\tfrac12 y^\top Q_i y+u_i^\top y,\\
 \partial_n\ell_i(o+y)&=b_i+v_i^\top y, & & y\in E,
 \end{aligned}
\]
where
\[
 \begin{aligned}
 Q_i&=\Pi P_i|_E, &
 u_i&=\Pi P_i(a_i-o), &
 c_i&=\ell_i(o),\\
 b_i&=n^\top P_i(a_i-o), &
 v_i&=-\Pi P_i n.
 \end{aligned}
\]
Group the indices with identical pairs $(Q_i,u_i)$ on this plane.

\begin{theorem}[Finite balance with component-specific precisions]
\label{thm:heterogeneous}
For a fixed affine readout as in Theorem~\ref{thm:flat}, preservation of
every strict decision at some $T>0$ is equivalent to preservation at
every $t\ge0$. Both are equivalent to the following scalar and vector
balances in every group $G$ on every facet hyperplane:
\begin{equation}
 \sum_{i\in G} w_i e^{c_i} b_i=0,\qquad
 \sum_{i\in G} w_i e^{c_i} v_i=0.
 \label{eq:heterogeneous-balance}
\end{equation}
\end{theorem}

\begin{proof}
The real-analytic potential $U$ tends to $-\infty$ as $\|x\|\to\infty$,
since there are finitely many positive definite precisions. The field
is a convex average of the vectors $P_i(a_i-x)$ and has at most linear
growth, so its flow is complete. If $\Phi_T$ preserves all strict
decisions, continuity maps the relative interior of each facet into
its supporting plane. Analyticity extends this inclusion to all of
$H$. Lemma~\ref{lem:sampled-tangency} then gives $F\cdot n=0$ on $H$.
Conversely, tangency on each supporting plane makes that plane invariant
by uniqueness. A trajectory from a strict cell cannot reach any of
its supporting planes in finite time by backward uniqueness.

It remains to characterize tangency. On $H$, multiplying
$F\cdot n$ by the positive mixture density gives
\[
 \sum_G A_G(y)e^{q_G(y)},\qquad
 A_G(y)=\sum_{i\in G}w_i e^{c_i}(b_i+v_i^\top y),
\]
where $q_G$ is the common tangential quadratic in group $G$.
If $\dim E=0$ there is one group and the vector balance is vacuous.
Otherwise choose a direction $\xi\in E$ such that
$\xi^\top(Q_G-Q_{G'})\xi\ne0$ whenever $Q_G\ne Q_{G'}$, and
$(u_G-u_{G'})^\top\xi\ne0$ whenever $Q_G=Q_{G'}$ and $G\ne G'$.
Each excluded condition is the zero set of a nonzero polynomial,
and their finite union cannot cover $E$.

For any $y_0\in E$, the restrictions $q_G(y_0+t\xi)$ have
distinct pairs of quadratic and linear coefficients in $t$.
Order them lexicographically by these coefficients. If the sum
vanishes identically, divide by the exponential of the largest
polynomial in this ordering and let $t\to+\infty$. Each other
exponential decays relative to the largest one faster than the growth
of its affine multiplier. The affine function $A_G(y_0+t\xi)$ for the
largest group must tend to zero, so it is identically zero on this
line. Repeating after its removal gives the same conclusion for
every group on the line. Since $y_0$ is arbitrary, all $A_G$ vanish
on $E$. This gives \eqref{eq:heterogeneous-balance}, which also
suffices for tangency.
\end{proof}

A singleton group contributes zero only when its center lies in $H$ and
$n$ is an eigenvector of its precision. Components with identical
tangential quadratics can cancel even when reflection does not
permute them. For example, take $H=\{x_1=0\}$ in $\mathbb R^2$ and
centers $(1,0)$ and $(-1,0)$ with precisions
$\operatorname{diag}(1,2)$ and $\operatorname{diag}(3,2)$,
respectively. Both have the same $Q_i$ and $u_i=0$, and both $v_i$
vanish. The weights $1$ and $e/3$ give the scalar balance
$e^{-1/2}-(e/3)\,3e^{-3/2}=0$. The partition into the two strict
halfspaces is invariant.
